\begin{frame}{Conclusões}
  \begin{itemize}
    \item A função de transferência descreve um filtro ativo passa-baixa de segunda ordem.
    \item A troca de $R_1$ e $R_3$ reduz a magnitude do ganho em DC de 10 para 1
      e altera o circuito de sobreamortecido para subamortecido.
    \item O circuito 1 tende a $-8\,\mathrm{V}$ para o degrau de $0{,}8\,\mathrm{V}$;
      o circuito 2 tende a $-5\,\mathrm{V}$ para o degrau de $5\,\mathrm{V}$.
    \item O LTspice e a bancada reproduzem a atenuação em altas frequências e
      as diferenças entre as configurações.
    \item O corte do circuito 1 foi registrado em $4{,}3\,\mathrm{Hz}$ e o pico do
      circuito 2 em $17{,}5\,\mathrm{Hz}$; os valores teóricos são $3{,}03$ e $17{,}92\,\mathrm{Hz}$.
  \end{itemize}

  \begin{alertblock}{Síntese da prática}
    A posição dos polos conecta o comportamento em frequência ao transitório.
    A comparação experimental considera a variação da tensão de entrada durante
    a varredura, utilizando o ganho $V_o/V_i$ em cada frequência.
  \end{alertblock}
\end{frame}
